%% kapa: exam=1 round=19 year=2008 period=1 form=A subject=civil_law subject_ko=민법 seq=21 number=21
%% source: Q-Net 감정평가사 기출문제 (cst00309) artlSeq=1001962 file=extracted/1008/1001962/2008년 제19회 감정평가사 1차시험 기출문제/1교시.hwp
%% conversion: kordoc 4.24.51 -> tools/gen1.py (auto); review=proofread-sonnet; flags=none
\begin{question}{21}
甲은 자신이 소유하고 있던 X토지를 乙에게 매도하였다. 乙은 매매대금을 모두 지급하고 X토지를 인도받아 점유․사용하고 있으나 乙 앞으로 소유권이전등기는 하지 않은 상태이다. 다음 설명 중 옳은 것은? (다툼이 있는 경우에는 판례에 의함)
\begin{choices}
  \item 미등기자인 乙이 丙에게 X토지를 매도하고 인도하였다면, 甲은 丙에 대하여 소유권에 기한 물권적 반환청구권을 행사할 수 있다.
  \item 乙이 X토지를 인도 받은 지 10년이 경과한 경우에는 甲은 乙에게 등기청구권의 소멸시효를 원용하면서 X토지의 반환을 청구할 수 있다.
  \item 제3자 丙이 X토지상에 무단으로 건물을 신축하였다면, 乙은 甲의 丙에 대한 물권적 청구권을 대위행사할 수 없다.
  \item 甲은 자신의 소유권에 기하여 乙에게 X토지에 대해서 소유물반환청구권을 행사할 수 없다.
  \item 제3자 丙이 무단으로 X토지에 건물을 신축하고 그 건물을 丁에게 임대하고 있다면, 甲은 丙에 대하여 건물철거청구를 할 수 없다.
\end{choices}
\end{question}
\answer{4}
